% GENERATED FILE. Edit canonical lessons, then run npm run generate:books.
% canonical-source-hash: fnv1a64:3850fd0678a30cf3
% canonical-lessons: ZH-C138-mei, ZH-W138-mei-every, ZH-C138-ziji, ZH-C138-bieren, ZH-C138-suoyou, ZH-C138-qita

\chapter{Every and Other}
\label{ch:zh-every-and-other}
\hlchaptermodality{138}

\begin{chapteropening}
\textbf{By the end of this chapter:} \emph{I can say every, myself, other people, all and the others.}

Complete the last lesson of chapter 138: I can say every, myself, other people, all and the others.
\end{chapteropening}

\section[\texorpdfstring{\zh{每} (měi)}{每 (měi)}]{\zh{每} (měi) — every; each}

\label{lesson:ZH-C138-mei}



\begin{quote}
Before the new one: say the Mandarin for to have time off, then the Mandarin for the holidays.
\end{quote}

\subsection*{You'll want to know: \zh{每}}
\textbf{\zh{每}} — \emph{měi} — \textquotedblleft{}every; each\textquotedblright{}.

\begin{quote}
\textbf{\zh{我们每个星期六都打扫}} \emph{wǒmen měi ge xīngqīliù dōu dǎsǎo} — We clean the house every Saturday.
\end{quote}

The new character, in pieces: \zh{每} — a short fall and a short bar on top, then a box: a vertical that turns right along the bottom, a bar that turns down the right side and hooks, a dot, a long bar across, and a dot. The next lesson puts it in your hand.

\begin{cousinweb}
\textbf{\zh{每}} (\emph{měi}) is inherited from Middle Chinese \emph{mwojX} (Guangyun, in Baxter's transcription), glossed 'frequently' and as a particle. 'Every' is not among the glosses, so how it got that sense is not settled.
\end{cousinweb}

\begin{grammarlens}[title={What you've built}]
One more word for pointing things out.
\end{grammarlens}

\subsection*{Your turn}
\begin{itemize}
\raggedright
  \item \emph{Say it:} \emph{měi}
  \item \emph{Say it:} \emph{měi}, once more
  \item \emph{Say it:} \emph{wǒmen měi ge xīngqīliù dōu dǎsǎo} — We clean the house every Saturday.
  \item \emph{Recall:} say the Mandarin for to have time off, then the Mandarin for the holidays, then say \emph{měi} again
\end{itemize}

\begin{tcolorbox}[breakable,colback=teal!4,colframe=teal!35!black,title={Before you move on}]
What does \zh{每} mean? (\textquotedblleft{}Every; each\textquotedblright{}.) Say it once more.
\end{tcolorbox}
\section[\texorpdfstring{\zh{每} (měi)}{每 (měi)}]{\zh{每} — a fall and a bar over a dotted box}

\label{lesson:ZH-W138-mei-every}



\begin{quote}
The word first, then the character inside it.

\begin{itemize}
\raggedright
  \item \emph{Recall:} say \emph{měi} — \textbf{R1}, one lesson back, before it decays
\end{itemize}

\begin{quote}
\zh{天} \zh{个}
\end{quote}

\emph{Read it:} each one aloud, with nothing said first — \emph{tiān}, \emph{gè}
\end{quote}

\subsection*{Script — a fall and a bar over a dotted box}
\hlblockfigure{\detokenize{figures/ZH-W138-mei-every-filmstrip.pdf}}{How \zh{每} is written, stroke by stroke}

\begin{quote}
\zh{每}
\end{quote}

\textbf{Two pieces, top to bottom.}

\textbf{The top, two strokes.} A short left-falling stroke, then a short horizontal running right from it.

\textbf{The box, five strokes.} A vertical that runs down the left and turns right along the bottom; a horizontal across the top that turns down the right side, past the bottom, and hooks to the left; a dot inside the upper half; a long horizontal across the middle, out past both sides; then a dot inside the lower half.

2 + 5 = seven.

\subsection*{Your turn}
\begin{itemize}
\raggedright
  \item \emph{Write it:} \zh{每}, all seven, with the model in front of you
  \item \emph{Cover:} write it again from memory, and count the pieces, not the strokes
\end{itemize}

\begin{tcolorbox}[breakable,colback=teal!4,colframe=teal!35!black,title={Before you move on}]
How many strokes in \zh{每}? (\textbf{Seven}.) Which stroke runs out past both sides? (\textbf{The long horizontal across the middle}.) Where are the two dots? (\textbf{Inside the box, one above that bar and one below it}.)

Source: \href{https://raw.githubusercontent.com/chanind/hanzi-writer-data/68d10a4b21150cae5e1ebbd223eed289cf32d90c/data/\%E6\%AF\%8F.json}{Hanzi Writer Data \zh{每}, pinned stroke paths and medians}.
\end{tcolorbox}
\section[zìjǐ]{zìjǐ — oneself}

\label{lesson:ZH-C138-ziji}



\begin{quote}
Before the new one: say the Mandarin for the holidays, then the Mandarin for each.
\end{quote}

\subsection*{You'll want to know: zìjǐ}
\emph{zìjǐ} — \textquotedblleft{}oneself\textquotedblright{}.

\begin{quote}
\emph{wǒ zìjǐ zuòfàn} — I cook for myself.
\end{quote}

Its characters come later; for now it is a word for the ear.

\begin{cousinweb}
\emph{zìjǐ} is built from \emph{zì} 'from' and \emph{jǐ} 'self', both inherited from Middle Chinese, \emph{dzijH} and \emph{kiX} (Guangyun, in Baxter's transcription), which glosses the second 'one's own body'.
\end{cousinweb}

\begin{grammarlens}[title={What you've built}]
One more word for pointing things out.
\end{grammarlens}

\subsection*{Your turn}
\begin{itemize}
\raggedright
  \item \emph{Say it:} \emph{zìjǐ}
  \item \emph{Say it:} \emph{zìjǐ}, once more
  \item \emph{Say it:} \emph{wǒ zìjǐ zuòfàn} — I cook for myself.
  \item \emph{Recall:} say the Mandarin for the holidays, then the Mandarin for each, then say \emph{zìjǐ} again
\end{itemize}

\begin{tcolorbox}[breakable,colback=teal!4,colframe=teal!35!black,title={Before you move on}]
What does zìjǐ mean? (\textquotedblleft{}Oneself\textquotedblright{}.) Say it once more.
\end{tcolorbox}
\section[biérén]{biérén — other people}

\label{lesson:ZH-C138-bieren}



\begin{quote}
Before the new one: say the Mandarin for each, then the Mandarin for oneself.
\end{quote}

\subsection*{You'll want to know: biérén}
\emph{biérén} — \textquotedblleft{}other people\textquotedblright{}.

\begin{quote}
\emph{wǒ de shū wǒ bù xiǎng gěi biérén} — I don't want to give my books to anyone else.
\end{quote}

Its characters come later; for now it is a word for the ear.

\begin{cousinweb}
\emph{biérén} is built from \emph{bié} 'different', as in \emph{biéde}, and \zh{人} 'person', both inherited from Middle Chinese, \emph{bjet} and \emph{nyin} (Guangyun, in Baxter's transcription).
\end{cousinweb}

\begin{grammarlens}[title={What you've built}]
One more word for pointing things out.
\end{grammarlens}

\subsection*{Your turn}
\begin{itemize}
\raggedright
  \item \emph{Say it:} \emph{biérén}
  \item \emph{Say it:} \emph{biérén}, once more
  \item \emph{Say it:} \emph{wǒ de shū wǒ bù xiǎng gěi biérén} — I don't want to give my books to anyone else.
  \item \emph{Recall:} say the Mandarin for each, then the Mandarin for oneself, then say \emph{biérén} again
\end{itemize}

\begin{tcolorbox}[breakable,colback=teal!4,colframe=teal!35!black,title={Before you move on}]
What does biérén mean? (\textquotedblleft{}Other people\textquotedblright{}.) Say it once more.
\end{tcolorbox}
\section[suǒyǒu]{suǒyǒu — all; every}

\label{lesson:ZH-C138-suoyou}



\begin{quote}
Before the new one: say the Mandarin for oneself, then the Mandarin for other people.
\end{quote}

\subsection*{You'll want to know: suǒyǒu}
\emph{suǒyǒu} — \textquotedblleft{}all; every\textquotedblright{}.

\begin{quote}
\emph{suǒyǒu de rén dōu lái le} — Everyone has come.
\end{quote}

Its characters come later; for now it is a word for the ear.

Two third tones: suǒyǒu is said suóyǒu, the first rising.

\begin{cousinweb}
\emph{suǒyǒu} is built from \emph{suǒ} 'a place', the \emph{suǒ} of \emph{suǒyǐ}, and \zh{有} 'to have', both inherited from Middle Chinese, \emph{srjoX} and \emph{hjuwX} (Guangyun, in Baxter's transcription). How the pair came to mean all is not settled.
\end{cousinweb}

\begin{grammarlens}[title={What you've built}]
One more word for pointing things out.
\end{grammarlens}

\subsection*{Your turn}
\begin{itemize}
\raggedright
  \item \emph{Say it:} \emph{suǒyǒu}
  \item \emph{Say it:} \emph{suǒyǒu}, once more
  \item \emph{Say it:} \emph{suǒyǒu de rén dōu lái le} — Everyone has come.
  \item \emph{Recall:} say the Mandarin for oneself, then the Mandarin for other people, then say \emph{suǒyǒu} again
\end{itemize}

\begin{tcolorbox}[breakable,colback=teal!4,colframe=teal!35!black,title={Before you move on}]
What does suǒyǒu mean? (\textquotedblleft{}All; every\textquotedblright{}.) Say it once more.
\end{tcolorbox}
\section[qítā]{qítā — other; the others}

\label{lesson:ZH-C138-qita}



\begin{quote}
Before the new one: say the Mandarin for other people, then the Mandarin for every single.
\end{quote}

\subsection*{You'll want to know: qítā}
\emph{qítā} — \textquotedblleft{}other; the others\textquotedblright{}.

\begin{quote}
\emph{wǒ xǐhuan zhège qítā de wǒ bú yào} — I like this one; I don't want the others.
\end{quote}

Its characters come later; for now it is a word for the ear.

\begin{cousinweb}
\emph{qítā} is built from \emph{qí}, which the dictionary lists as a particle, and \zh{他}, glossed 'not I', both inherited from Middle Chinese, \emph{gi} and \emph{tha} (Guangyun, in Baxter's transcription).
\end{cousinweb}

\begin{grammarlens}[title={What you've built}]
One more word for pointing things out.
\end{grammarlens}

\subsection*{Your turn}
\begin{itemize}
\raggedright
  \item \emph{Say it:} \emph{qítā}
  \item \emph{Say it:} \emph{qítā}, once more
  \item \emph{Say it:} \emph{wǒ xǐhuan zhège qítā de wǒ bú yào} — I like this one; I don't want the others.
  \item \emph{Recall:} say the Mandarin for other people, then the Mandarin for every single, then say \emph{qítā} again
\end{itemize}

\begin{tcolorbox}[breakable,colback=teal!4,colframe=teal!35!black,title={Before you move on}]
What does qítā mean? (\textquotedblleft{}Other; the others\textquotedblright{}.) Say it once more.
\end{tcolorbox}
